% ============================================================
%  ch15  第五讲 §5.5 自举（bootstrap）
%  原文行 3075–3210
% ============================================================

\subsection{自举}

要把几何 Satake 对应用于（几何）朗兰兹纲领，重要的是不假设 $k$ 代数闭，并把它推广为 $\Sat_{X^I}$（对各个 $I$）与涉及 Langlands 对偶群 $\hat G$ 的某些范畴之间的等价。本小节说明如何自举先前的结果以得到这样的推广。

先讲一些一般性的形式框架。

\begin{dfn}
设 $K$ 是域 $E$ 上的线性代数群，$\mC$ 是 $E$-线性 Abel 范畴且容许无穷直和（于是可以与 $E$-向量空间张量）。$\mC$ 中 $K$ 的\textbf{表示}是一个对象 $M$ 连同态射 $\rho: M\to M\otimes_{E}\mO_K$，满足余作用通常的交换图。$K$-表示的态射 $f: (M_1,\rho_1)\to (M_2,\rho_2)$ 是 $\mC$ 中与 $\rho_i$ 相容的态射 $f:M_1\to M_2$。$\mC$ 中 $K$-表示的范畴记作 $\on{Rep}(K,\mC)$。

一般地，若 $\mC$ 未必容许无穷直和，先过渡到其归纳完备化 $\on{Ind}\mC$，再把 $\on{Rep}(K,\mC)$ 定义为 $\on{Rep}(K,{\on{Ind}\mC})$ 中那些底对象 $M$ 属于 $\mC$ 的 $(M,\rho)$ 构成的全子范畴。容易验证：当 $\mC$ 容许无穷直和时，新定义与旧定义一致。
\end{dfn}

\begin{ex}
若 $\mC=\on{Vect}_E$ 是有限维 $E$-向量空间范畴，则 $\on{Rep}(K,\mC)$ 就是 $K$ 的有限维代数表示范畴，为简洁记作 $\Rep_E(K)$。
\end{ex}

\begin{ex}\label{KLoc}
设 $\Ga$ 是投射有限群，$\on{Rep}^c_\Ql(\Ga)$ 记 $\Ga$ 的连续 $\ell$-进表示范畴。设 $\on{Rep}(K,{\on{Rep}^c_\Ql(\Ga)})$ 是上述 $K$ 在 $\on{Rep}^c_\Ql(\Ga)$ 中表示的范畴。则容易看出，$\on{Rep}(K,{\on{Rep}^c_\Ql(\Ga)})$ 的对象恰是 $\Ga\times K$ 的表示，它对第一个因子连续、对第二个因子代数。
\end{ex}

注意若 $F:\mC\to \mD$ 是加性函子，它典范地提升为函子 $F:\on{Rep}(K,{\mC})\to \on{Rep}(K,{\mD})$。

表示论中若干经典构造可推广到此框架，择要列举。给定 $K$-线性表示 $V$ 与 $M\in \on{Rep}(K,{\mC})$，张量积 $M\otimes V$ 作为 $\on{Rep}(K,{\mC})$ 的对象有意义。若 $\mC$ 是幺半 Abel 范畴且 $M,N\in \on{Rep}(K,{\mC})$，则 $M\otimes N$ 带有自然的对角 $K$-作用，因此 $\mC$ 的幺半结构可升级为 $\on{Rep}(K,{\mC})$ 上的幺半结构。还可定义不变量 $M^K$ 为映射 $\rho:M\to M\otimes \mO_K$ 与映射 $M=M\otimes_EE\stackrel{e_K}{\to} M\otimes_E\mO_K$ 的等化子。则 $M\mapsto M^K$ 是从 $\on{Rep}(K,{\mC})$ 到 $\mC$ 的左正合函子。

现在设 $k$ 是未必代数闭的域，$\bar k$ 是其可分闭包，$\Ga_k=\Gal(\bar k/k)$。设 $G$ 是 $k$ 上连通约化群，$\underline G=G\otimes_k\mO$，$\Gr$ 记 $\underline G$ 的仿射 Grassmann 簇。设 $\hat G=\Aut^{\otimes}(\on{H}^*)$ 是 $\Sat_{G_{\bar k}}$ 的 Tannaka 群，配以命题 \ref{canonical pinning} 的典范带钉\footnote{要把它与通常组合构造给出的群等同，其实并不需要定理 \ref{birth of dual group}。}。下面先说明如何把完整 Langlands 对偶群 ${^L}G$ 定义为 $\hat G$ 与 $\Ga_k$ 的半直积。事实上将看到 ${^L}G$ 有两个版本。

如 \cite{RZ} 所解释，有 $\Ga_k$ 在 $\hat G$ 上的典范作用 $\on{act}^{\on{geo}}$，定义如下。每个 $\ga\in \Ga_k$ 诱导自同构 $\Gr\otimes \bar k\cong \Gr\otimes \bar k$，故拉回 $\ga^*$ 诱导 $\Sat_{G_{\bar k}}$ 的张量自同构。此外有张量函子的典范同构 $\al_\ga: \on{H}^*\cong \on{H}^*\circ \ga^*$。于是对 $g\in \hat G$（看作张量函子 $\on{H}^*$ 的自同构），$\ga(g)$ 定义为 $\on{H}^*$ 的张量自同构
\begin{equation}\label{gaaction}
\on{H}^*\stackrel{\al_\ga}{\longrightarrow}\on{H}^*\circ\ga^*\stackrel{g\circ \id}{\to}\on{H}^*\circ\ga^*\stackrel{\al^{-1}_\ga}{\to}\on{H}^*.
\end{equation}
由于 $\al_\ga$ 保持上同调次数，$\on{act}^{\on{geo}}$ 保持 $2\rho:\bG_m\to \hat G$，从而由推论 \ref{canonical TB} 保持 $\tilde T\subset \tilde B\subset \hat G$；但它不保持命题 \ref{canonical pinning} 中的主幂幺元 $N$。

设 $\hat G_\ad$ 是 $\hat G$ 的伴随商。则余特征 $\bG_m\stackrel{2\rho}{\to} \hat G\to \hat G_\ad$ 有唯一的平方根，记作 $\rho_\ad$。设 $\on{cycl}:\Ga_k\to\bZ_\ell^\times$ 是由 $\Ga_k$ 在 $\bar k$ 的 $\ell^\infty$-单位根上的作用定义的分圆特征。定义同态
\[\chi:\Ga_k\stackrel{\on{cycl}}{\to} \bZ_\ell^\times\stackrel{\rho_\ad}{\to} \hat G_{\ad}(\Ql),\]
它通过内自同构诱导 $\Ga_k$ 在 $\hat G$ 上的作用 $\Ad_{\chi}:\Ga_k\to\Aut(\hat G)$。

\begin{lem}\label{comparison}
作用 $\on{act}^{\on{alg}}:=\on{act}^{\on{geo}}\circ \Ad^{-1}_{\chi}$ 保持命题 \ref{canonical pinning} 中定义的 $\hat G$ 的带钉。
\end{lem}

\begin{proof}
这由第一 Chern 类 $c_1(\mO(1))$ 属于 $\on{H}^2(\Gr_G\otimes\bar k, \Ql(1))^{\Ga_k}$ 推出。细节见 \cite[Proposition A.6]{RZ}。
\end{proof}

\begin{rmk}
本注后文不需要。回忆总有 $\Ga_k$ 通过外自同构在 $G\otimes \bar k$ 上的作用 $\psi: \Ga_k\to \on{Out}(G\otimes \bar k)$。另一方面，由定理 \ref{birth of dual group} 有典范同构 $\on{Out}(G\otimes\bar k)\cong \on{Out}(\hat G)=\Aut(\hat G, \tilde B, \tilde T, N)$。如 \cite{RZ} 所解释，在此同构下 $\on{act}^{\on{alg}}$ 与 $\psi$ 一致。
\end{rmk}

把 $G$ 的完整 Langlands 对偶群定义为
\[{^L}G^{\on{alg}}:= \hat G\rtimes_{\on{act}^{\on{alg}}}\Ga_k.\]
因 $\Ga_k$ 投射有限，且其在 $\hat G$ 上的作用 $\on{act}^{\on{alg}}$ 通过有限商分解，可把 ${^L}G^{\on{alg}}$ 看作 pro-代数群。

另一方面，按上述讨论，考虑
\[{^L}G^{\on{geo}}:= \hat G\rtimes_{\on{act}^{\on{geo}}}\Ga_k\]
也很自然，它与 $\Sat_G$ 联系更紧密。把 ${^L}G^{\on{geo}}$ 看作拓扑群。

作为一条勘误，我们指出 \cite[Corollary A.8]{RZ} 有误，正确陈述如下。

\begin{lem}\label{geom Langlands dual}
设 $\chi: \Ga_k\to \hat G_\ad(\Ql)$ 可提升为连续同态 $\tilde \chi:\Ga_k\to \hat G(\Ql)$。则有（拓扑群的）自然同构
\[{^L}G^{\on{geo}}\cong{^L}G^{\on{alg}},\quad (g,\ga)\mapsto (g\tilde\chi(\ga),\ga).\]
\end{lem}

\begin{rmk}\label{lifting condition}
提升 $\tilde \chi$ 在下列情形存在。
\begin{enumerate}
\item[(i)] 余特征 $\rho_\ad$ 提升为余特征 $\tilde \rho: \bG_m\to \hat G$。例如 $\hat G=\GL_n$ 时成立。注意并不要求 $2\tilde \rho=2\rho$。
\item[(ii)] 分圆特征 $\on{cycl}: \Ga_k\to \bZ_\ell^\times$ 有平方根，即存在半 Tate 扭 $\Ql(\frac{1}{2})$。例如 $\on{char} k=p>0$，或 $k$ 是 $p$-进域（$p\neq \ell$）时成立。
\end{enumerate}
\end{rmk}

\begin{ex}\label{geom v.s. alg dual}
这是两个群不同构的例子。取 $k=\bR$、$G=\on{PGL}_2(\bR)$。则 ${^L}G^{\on{alg}}=\SL_2\times \Gal(\bC/\bR)$，而 ${^L}G^{\on{geo}}=\SL_2\rtimes \Gal(\bC/\bR)$，其中复共轭通过
\[\begin{pmatrix} a & b\\ c & d \end{pmatrix}\mapsto \begin{pmatrix} a & -b\\ -c & d \end{pmatrix}\]
作用在 $\SL_2$ 上。注意 ${^L}G^{\on{alg}}$ 的中心同构于 $\bZ/2\times \bZ/2$，而 ${^L}G^{\on{geo}}$ 的中心同构于 $\bZ/4$，故两者不同构。
\end{ex}

虽然 ${^L}G^{\on{geo}}$ 不能看作 pro-代数群，但总可把它嵌入更大的 pro-代数群，如下。令
\begin{equation}\label{Deligne modification}
\hat G^T:=\hat G\rtimes\bG_m,
\end{equation}
其中 $\bG_m$ 通过 $\rho_\ad:\bG_m\to \hat G_\ad$ 以内自同构作用在 $\hat G$ 上。例如若 $\hat G=\SL_2$，则 $\hat G^T=\GL_2$。作用 $\on{act}^{\on{alg}}$ 自然延拓为 $\Ga_k$ 在 $\hat G^T$ 上的作用，仍记作同一记号。令
\[{^L} G^T=\hat G^T\rtimes_{\on{act}^{\on{alg}}}\Ga_k,\]
它可看作 pro-代数群。以 $d: \hat G^T\to \bG_m$ 记自然投影，则有短正合列
\[1\to {^L}G^{\on{geo}}\to {^L}G^T\stackrel{d\times \on{cycl}^{-1}}{\longrightarrow}\bG_m\to 1.\]

\begin{rmk}
注意 $\hat G^T$ 其实是 $G$ 被 $\bG_m$ 的中心扩张的对偶群。看来 ${^L}G^T$ 正是 \cite[Definition 5.3.2]{BuGe} 中出现的 $C$-群。如上面所提示的，也如 \cite{BuGe} 所建议，为了给自守表示配 Galois 表示，这似乎是更自然的对象。
\end{rmk}

以 $\on{Rep}^c_{\Ql}({^L}G^{\on{geo}})$ 记 ${^L}G^{\on{geo}}$ 的有限维连续表示中限制到 $\hat G$ 为代数者的范畴。类似地定义 $\on{Rep}^c_{\Ql}({^L}G^T)$。设 $B$ 是 $\bG_m$ 的正则表示，经同态 $d\times \on{cycl}^{-1}$ 看作 $\on{Rep}_{\Ql}({^L}G^T)$ 的归纳完备化中的代数对象。回忆对幺半范畴 $\mC$ 与代数对象 $B\in \mC$，可谈 $\mC$ 中的 $B$-模，其范畴记作 $B\on{-mod}(\mC)$。则有等价
\[B\on{-mod}(\on{Rep}^c_{\Ql}({^L}G^T))\simeq \on{Rep}^c_{\Ql}({^L}G^{\on{geo}}),\quad V\mapsto V\otimes_B\Ql,\]
其中 $B\to \Ql$ 是余单位映射（对偶于 $1\in \bG_m$）。

现在以 $\Sat_G$ 记 $\Gr$ 上 $L^+G$-等变反常层范畴，配以卷积积（注意 \eqref{Lus conv} 定义在 $k$ 上）。例如若 $G=\{e\}$ 平凡，这就是 $\Ga_k$ 的连续 $\ell$-进表示范畴 $\on{Rep}^c_{\Ql}(\Ga_k)$，配以通常张量积。下面是几何 Satake 等价的完整版本。

\begin{thm}\label{full Sat equiv}
超上同调函子 $\on{H}^*:\Sat_G\to \on{Rep}^c_{\Ql}(\Ga_k)$ 提升为典范等价
\begin{equation}\label{full Sat}
F:\Sat_G\simeq \on{Rep}^c_{\Ql}({^L}G^{\on{geo}}),
\end{equation}
且下图交换（差一个典范同构）
\[\begin{CD}
\Sat_G@>\simeq>> \on{Rep}^c_{\Ql}({^L}G^{\on{geo}})\\
@V{\text{拉回}}VV@VV{\text{限制}}V\\
\Sat_{G_{\bar k}}@>\simeq>> \on{Rep}_{\Ql}(\hat G).
\end{CD}\]
\end{thm}

\begin{proof}
勾勒证明（另一途径见 \cite[\S~5]{Ri}）。因 $\Gr$ 上的层拉回到 $\Gr\otimes \bar k$ 是 $\Ga_k$-等变的，由一般形式框架（\cite[Lemma A.3]{RZ}）可知超上同调按定量所述分解，且图交换。剩下要证它是等价，将构造拟逆函子来证。

为说明想法，假设 $G$ 分裂，从而每个 $L^+G$-轨道定义在 $k$ 上。若 $G$ 非分裂，可先过渡到使 $G_{k'}$ 分裂的有限扩张 $k'/k$，再做进一步的 Galois 下降。

对 Schubert 簇 $\Gr_{\leq \mu}$，以 $\IC_\mu$ 记其相交上同调复形（故 $\IC_\mu|_{\Gr_\mu}=\Ql[(2\rho,\mu)]$）。先考虑 $\Sat_G$ 中由 $\mu\in \xcoch(T)^+$、$i\in \bZ$ 的 $\IC_\mu(i)$ 的直和张成的全子范畴 $\Sat_G^T$。

\begin{lem}~
\begin{enumerate}
\item 范畴 $\Sat_G^T$ 是 Tannaka 范畴，纤维函子由 $\on{H}^*$ 给出。
\item Tannaka 群 $\Aut^{\otimes}(\on{H}^*)$ 同构于 \eqref{Deligne modification} 中定义的 $\hat G^T$。
\end{enumerate}
\end{lem}

\begin{proof}
(1) 事实上由推论 \ref{Satake fiber}，$\Sat_G^T$ 对卷积积封闭。又因命题 \ref{existence of comm} 中的交换约束 $c$ 是 $\Ga_k$-等变的（因 $c_{\on{vect}}$ 是 $\Ga_k$-等变的），$\Sat_G^T$ 对称幺半。于是如定理 \ref{Tannakian structure} 的论证，它是 Tannaka 范畴。

(2) 记该 Tannaka 群为 $\tilde G^T$。有正合列
\[1\to \hat G\to \tilde G^T\stackrel{d}{\to} \bG_m\to 1,\]
其中第一个映射由拉回 $\Sat_G^T\to \Sat_{G_{\bar k}}$ 诱导，第二个由 $\IC_e(i)$（$i\in\bZ$）的直和张成的全子范畴给出。投影 $d:\tilde G^T\to \bG_m$ 有分裂 $\bG_m\to \tilde G^T$，对应于 $\on{H}^*(\mA)$ 上的下述分次
\[\on{H}^*(\mA)=\bigoplus_i \Ql(i)\otimes \Hom(\Ql(i), \on{H}^*(\mA)).\]
于是如引理 \ref{geom Langlands dual}，用 $c_1(\mO(1))\in \on{H}^2(\Gr_G\otimes\bar k,\bQ_\ell(1))^{\Ga_k}$ 即得 $\tilde G^T=\hat G^T$。
\end{proof}

固定 $\on{H}^*$ 的拟逆 $S^T: \Rep_\Ql(\hat G^T)\simeq \Sat_G^T$。

现在考虑整个范畴 $\Sat_G$，将用到本节开头的一般形式框架。设 $R^T$ 是 $\hat G^T$ 上正则函数环，经左正则表示看作 $\on{Ind}\on{Rep}_{\Ql}(\hat G^T)$ 的对象。令 $\mR^T=S^T(R^T)\in \on{Ind}\Sat^T_G$ 为对应对象。因 $R^T$ 还有右正则表示，$\mR^T$ 其实是 $\on{Rep}(\hat G^T,{\on{Ind}\Sat_G})$ 中的对象。因 $G$ 分裂，$\Rep^c_\Ql({^L}G^T)=\on{Rep}(\hat G^T,{\on{Rep}^c_{\Ql}(\Ga_k)})$。$\on{Rep}(\hat G^T,{\on{Rep}^c_{\Ql}(\Ga_k)})$ 中的对象 $V$ 可看作带 $\hat G^T$ 作用的 $\Spec k$ 上层。于是由一般形式框架，$e_*V\star\mR^T$ 作为 $\Rep(\hat G^T, \on{Ind}\Sat_G)$ 的对象有意义，故 $(e_*V\star \mR^T)^{\hat G^T}$ 是 $\on{Ind}\Sat_G$ 的对象。由此得到张量函子
\[\on{Rep}(\hat G^T,{\on{Rep}^c_{\Ql}(\Ga_k)})\to \on{Ind}\Sat_G,\quad V\mapsto (e_*V\star \mR^T)^{\hat G^T},\]
仍记作 $S^T$，因为它在 $\Rep_\Ql(\hat G^T)\subset \on{Rep}(\hat G^T,{\on{Rep}^c_{\Ql}(\Ga_k)})$ 上的限制恰是开头固定的函子 $S^T$。

回忆把 $B=\mO_{\bG_m}$ 看作 $\on{Rep}(\hat G^T,{\on{Rep}^c_{\Ql}(\Ga_k)})$ 的对象。则
\[\mB:=S^T(B)= \IC_e\otimes B.\]
因此 $S^T$ 把 $\on{Rep}(\hat G^T,{\on{Rep}^c_{\Ql}(\Ga_k)})$ 中的 $B$-模对象送到 $\on{Ind}\Sat_G$ 中的 $\mB$-模对象。于是可定义
\[S: \on{Rep}^c_{\Ql}({^L}G^{\on{geo}})\simeq B\on{-mod}(\on{Rep}^c_{\Ql}({^L}G^T)) \to \on{Ind}\Sat_G:\quad V\mapsto S^T(V)\otimes_{B}\Ql,\]
其中如前 $B\to\Ql$ 是余单位映射。

注意对 $V\in \on{Rep}(\hat G^T,{\on{Rep}^c_{\Ql}(\Ga_k)})$，
\[\on{H}^*\circ S^T(V)=\on{H}^*((e_*V\star \mR^T)^{\hat G^T})=(V\otimes R^T)^{\hat G^T}=V,\]
看作 ${^L}G^{\on{geo}}$ 的表示。由此可知 $S$ 正是 $\on{H}^*:\Sat_G\to \on{Rep}^c_{\Ql}({^L}G^{\on{geo}})$ 所需的拟逆。细节留作练习。
\end{proof}

\begin{rmk}\label{simplify}
若注 \ref{lifting condition} 中某一条件成立，定理的证明可简化。例如若 $\rho\in\xch(T)$ 或存在半 Tate 扭，可定义标准化的相交上同调复形
\[\IC_\mu^{\on{norm}}=\IC_\mu((2\rho,\mu)),\]
并定义由这些标准化 IC 层的直和张成的子范畴 $\Sat_G^N\subset \Sat_G$。由推论 \ref{Satake fiber}，它是 $\Sat_G$ 的幺半子范畴，且拉回诱导幺半范畴的等价
\[\Sat_G^N\simeq \Sat_{G_{\bar k}}.\]
因此 $\hat G$ 上正则函数环 $R$ 给出对象 $\mR\in \on{Rep}(\hat G,{\on{Ind}\Sat_G})$，可用它代替 $\mR^T$ 直接定义函子 $S$。
\end{rmk}

因 ${^L}G^{\on{alg}}$ 可看作 pro-代数群，有代数表示范畴 $\on{Rep}_\Ql({^L}G^{\on{alg}})$。若提升 $\tilde \chi:\Ga_k\to \hat G(\Ql)$ 存在，则经引理 \ref{geom Langlands dual} 中的自然同构，可把 $\on{Rep}_\Ql({^L}G^{\on{alg}})$ 看作 $\on{Rep}^c_{\Ql}({^L}G^{\on{geo}})$ 的全子范畴。在 $\rho\in \xch(T)$ 或存在半 Tate 扭的假设下，\cite{RZ} 说明了如何在等价 \eqref{full Sat} 下选出 $\Sat_G$ 中对应于 $\on{Rep}_\Ql({^L}G^{\on{alg}})$ 的全子范畴（\emph{同上}记作 $\on{P}_{L^+G}^f(\Gr_G)$）。当 $k$ 有限时，陈述尤为简洁。

\begin{prop}
设 $k=\bF_q$ 有限。以 $\Sat_G^0\subset \Sat_G$ 记权为零的半单纯反常层构成的全子范畴。则等价 \eqref{full Sat} 限制为等价
\[F:\Sat_G^0\simeq \on{Rep}_\Ql({^L}G^{\on{alg}}).\]
\end{prop}

\begin{proof}
见 \cite[Remark A.9, Proposition A.10]{RZ}。注意对 $\mA\in\Sat_G^0$，$\Ga_k$ 在 $F(\mA)$ 上的作用与 $\Ga_k$ 在 $\oplus_i \on{H}^i(\mA(\frac{i}{2}))$ 上的自然作用一致。
\end{proof}

下面描述对所有 $\Sat_{X^I}$ 的几何 Satake 等价的一个版本。为简单起见，假设 $G$ 分裂，且 $\rho\in \xch(T)$ 或存在半 Tate 扭。则如上所述，有完全忠实函子
\[S: \Rep(\hat G)\simeq \Sat_G^N\to \Sat_G.\]
设 $X$ 是 $k$ 上的曲线。对有限非空集 $I$，以 $\Loc(X^I)$ 记 $X^I$ 上 $\Ql$-局部系范畴，另有 $\Rep(\hat G^I,\Loc(X^I))$。联系这些范畴的有下述函子。设 $\phi:J\twoheadrightarrow I$。
\begin{enumerate}
\item[(i)] 限制 $\Delta(\phi)^*: \on{Rep}(\hat G^J,{\on{Loc}(X^J)})\to \on{Rep}(\hat G^I, {\on{Loc}(X^I)})$，其中 $\hat G^I$ 经嵌入 $\hat G^I\to \hat G^J$ 作用在 $\Delta(\phi)^*\mF$ 上。
\item[(ii)] 外积 $\boxtimes_{i\in I}:\prod_{i\in I} \on{Rep}(\hat G^{J_i},{\on{Loc}(X^{J_i})})\to \on{Rep}(\hat G^J,\on{Loc}(X^J))$。
\end{enumerate}

\begin{thm}\label{multiple point Sat}
存在典范等价
\[S_I: \on{Rep}(\hat G^I, \on{Loc}(X^I))\simeq \Sat_{X^I},\]
它与限制函子及外积函子在显然的意义下相容。
\end{thm}

\begin{rmk}
本定理的一个版本最早由 Gaitsgory 证明（\cite{Ga2}）。事实上他在两边都考虑了更大的范畴：左端为 $\on{Rep}(\hat G^I, \on{P}(X^I))$，右端为 $\on{P}_{(L^+G)_{X^I}}(\Gr_{X^I})$。然而如 (1) 中那样的限制函子在这些更大范畴上没有定义（见注 \ref{triv group}）。我们的途径与\emph{同上}不同。
\end{rmk}

下述推论在 \cite{Laf} 中起关键作用。

\begin{cor}
对每个 $I$ 有完全忠实函子 $S_I: \on{Rep}(\hat G^I)\to \Sat_{X^I}$，且对每个满射 $\phi: J\twoheadrightarrow I$ 有函子的典范同构 $\Delta(\phi)^*S_J\cong S_I\circ \Res_{\hat G^I}^{\hat G^J}$，其中 $\Res_{\hat G^I}^{\hat G^J}$ 是嵌入 $\hat G^I\to \hat G^J$ 对应的通常限制函子。
\end{cor}

说明该定理的证明思路。首先断言函子 $\mH^* :\Sat_{X^I}\to \on{Loc}(X^I)$ 可典范地提升为函子
\[F_I: \Sat_{X^I}\to \on{Rep}(\hat G^I,{\on{Loc}(X^I)}).\]
即：设 $\eta^I$ 是 $X^I$ 的通用点，$\overline{\eta^I}$ 是 $\eta^I$ 上的几何点。由于限制函子
\[i_{\overline{\eta^I}}^*: \on{Loc}(X^I)\to \on{Rep}^c_{\Ql}(\Gal(\overline{\eta^I}/\eta^I))\]
完全忠实，只需说明每个 $\mH^*(\mA)_{\overline{\eta^I}}$ 带有 $\hat G^I$ 的典范作用。然而 $(\Gr_{X^I})_{\eta^I}\cong (\prod_i \Gr_X)_{\eta^I}$，故所需提升由定理 \ref{full Sat equiv} 推出。

为反方向定义函子，需要 $\on{Rep}(\hat G^I, {\on{Ind}\Sat_{X^I}})$ 中类似于 $\Gr$ 上层 $\mR^T$（或在我们的假设下为注 \ref{simplify} 中的层 $\mR$）的对象。定义
\[\mR_X:= S(R)_X,\]
它（由定理 \ref{full Sat equiv} 中的论证）是 $\on{Rep}(\hat G, {\on{Ind}\Sat_X})$ 的对象。对一般 $I$，定义
\[\mR_{X^I}:=\boxstar_{i\in I} \mR_X\]
为 $\mR_X$ 的自外卷积积，自然看作 $\on{Rep}(\hat G^I,{\on{Ind}\Sat_{X^I}})$ 的对象。

于是可定义函子 $S_{I}: \on{Rep}(\hat G^I,{\on{Ind}\on{Loc}(X^I)})\to \on{Ind}\Sat_{X^I}$ 为
\[S_{I}(\mF)= ((e_I)_*\mF[|I|]\ostar\mR_{X^I})^{\hat G^I}.\]
断言：对每个 $\mA\in\Sat_{X^I}$ 有典范同构 $S_IF_I(\mA)\cong \mA$；对每个 $\mF\in \on{Rep}(\hat G^I, {\on{Loc}(X^I)})$ 有典范同构 $F_IS_I(\mF)\cong \mF$。注意由 ULA 性质，限制函子
\[\Sat_{X^I}\to \on{P}_{((L^+G)_{X^I})_{\eta^I}}((\Gr_{X^I})_{\eta^I})\]
完全忠实。因此只需在 $\eta^I$ 上构造这些同构，这就化归到定理 \ref{full Sat equiv}。

由 $\mR_{X^I}$ 的分解性质容易验证 $S_I$ 与外卷积积相容；由卷积积作为融合积的解释，$F_I$ 与到对角线的限制相容。定理得证。
